\chapter{實驗設計與結果}\label{chap:results}
\begingroup
\setlength{\emergencystretch}{2em}

\section{統計分析資料與比較範圍}

本章依第三章的兩項實驗安排，先呈現實驗一的 Baseline 與 Old／New 集成結果，再比較實驗二的特徵選取條件。歷史 Baseline 與固定年份集成使用不同驗證制度，分開彙整；公司群聚配對檢定只用於同一固定年份協定下的方法。集成與 FS 比較使用同一份資料切分：Old 擬合期間為 1999--2011 年，New 擬合期間為 2012--2013 年，2014 年分為 DSEL 與 Calibration，2015--2018 年為 Test。Test 含 12,282 筆公司年度紀錄，其中 287 筆屬破產正類，來自 3,700 家公司，其中 119 家為破產公司。

固定年份集成與 FS 共使用十個種子與九種特徵設定，合計 90 組設定。進行統計分析前，先核對各組測試紀錄與標籤，再以保存的機率及門檻重算 AUC、F1-score 與 Recall；所比較的 560 個輸出均與原紀錄一致。統計分析沿用原有模型、門檻與靜態選模結果。

表中 pooled 指標先於每個種子內合併四個測試年度計算，再跨十個種子彙整；種子標準差描述程序隨機變異。差值的信賴區間與近似檢定則使用第~\ref{sec:method-statistics}~節的公司群聚 bootstrap，將同公司跨年的紀錄一併重抽。兩項實驗共 165 項比較共同採 Holm 校正，與只將十個種子視為獨立樣本的檢定具有不同推論單位。

Baseline 的 Old／New 使用 15 組年份邊界，每組包含四種採樣；Retrain 的每種採樣只有一筆完整期間結果。其表格不加上跨種子的標準差，也不把重疊邊界當成獨立重複。與固定年份集成尚未完成同協定的重新調參比較，因此兩類結果的分數差不作為集成效益之檢定證據。

\section{實驗一：Baseline 分類器比較}\label{sec:results-baseline}

表~\ref{tab:results-baseline} 彙整七組已保存的分類器／版本紀錄。Old 與 New 先在同一邊界內平均四種採樣，再平均 15 組邊界；Retrain 直接平均四種採樣。這個彙整方式讓每個邊界與每種採樣具有相同權重，沒有依 Test 挑出最佳組合。各採樣的完整結果列於附錄~\ref{sec:results-baseline-full}。

\input{generated/chapter-results-20261027/baseline-summary.tex}

七組紀錄的 New 平均 AUC 與 F1-score 均高於對應 Old。例如 XGBoost 擴充版本的 Old／New AUC 為 0.7091／0.8589，F1-score 為 0.1138／0.2475；其 Retrain AUC 為 0.8386。這些數值反映本資料中近期期間模型的描述性表現，並未排除樣本量、驗證來源或歷史執行設定差異的影響。

Recall 的方向並不完全相同。Logistic Regression 的 Old Recall 為 0.6305，高於 New 的 0.3710，但 Old F1-score 僅為 0.0691，低於 New 的 0.1493。因此，較高 Recall 不能單獨解釋為較好的分類表現，也需要同時檢視正類預測的精確程度。

XGBoost 的兩種版本分別保留，TabICL 則按原結果標記呈現；未保存的後端與驗證資訊不從 Test 指標反推。這些紀錄可檢查資料期間與採樣的敏感度，尚不足以建立同等調參預算下的演算法排名。後續集成採用第三章所列的固定 XGBoost。

\section{實驗一：Old／New 集成比較}\label{sec:results-experiment1}

\subsection{六個基礎模型與靜態成員數}

本部分使用完整 18 個特徵與固定年份配置。表~\ref{tab:results-pool-models} 先列六個池內模型的測試結果，使後續整合結果可連同成員本身的表現解讀。

\input{generated/chapter-results-20261027/pool-models.tex}

例如 New-under 的 AUC、F1-score 與 Recall 為 0.8067、0.2332 與 0.3206，顯示模型池內各成員的表現並不相同。這項個別模型結果未納入原先相對 Equal6 的 165 項檢定，故不補作其與 V-static 或動態方法之顯著性主張。

\input{generated/chapter-results-20261027/static-size.tex}

表~\ref{tab:results-static-size} 將 49 個混合候選依成員數分組。兩模型候選的平均 AUC 為 0.7701，六模型為 0.7815；F1-score 並非在每次增加成員時都上升。這些平均值涵蓋不同模型組合，不能解釋為只改變成員數所造成的效果。實際採用的 V-static 仍由 Calibration AP 決定，不依本表的 Test 平均值選擇成員數。

\subsection{集成方法與配對檢定}

先比較使用完整 18 個特徵的集成結果。Old3 與 New3 各整合同一期間的三種採樣模型，Equal6 使用全部六模型；V-static 則依各種子的 Calibration AP，從 49 個混合子集中選定成員。動態方法使用相同六模型池與 DSEL，結果見表~\ref{tab:stats-exp1-descriptive}。

\input{generated/statistics-20261027T051346477136Z/experiment1-descriptive.tex}

各方法相對 Equal6 的 AUC 差值、公司群聚信賴區間與 Holm 校正值列於表~\ref{tab:stats-exp1-auc}。正差值代表相對基準的 AUC 增加，負差值代表降低。區間為逐項 95\% 區間，顯著性則依全體比較校正後的 \(p_H\) 判斷。

\input{generated/statistics-20261027T051346477136Z/experiment1-auc-tests.tex}

V-static 的 AUC 與 F1-score 相對 Equal6 分別增加 0.0337 與 0.0540，Holm 校正後 \(p_H\) 分別為 0.00632 與 0.02416。AUC 差值的逐項 95\% 信賴區間為 [0.0198, 0.0480]。其靜態成員由 Calibration 決定，因此該差異屬於驗證選模程序與完整六模型等權平均之比較。相對地，Old3 的 AUC 與 F1-score 差值分別為 \(-0.0681\) 與 \(-0.0647\)，兩者 \(p_H\) 均為 0.00330。

DCS-OLA 與 DCS-LCA 的 AUC 相對 Equal6 分別降低 0.0430 與 0.0446，兩者 \(p_H=0.01860\)。New3、KNORA-E 與 KNORA-U 的 AUC 比較未通過整體 Holm 校正；所有方法相對 Equal6 的 Recall 比較亦未達校正後顯著。這些結果顯示，不能預設動態選擇或跨期間整合必然優於所有控制組。

F1-score 與 Recall 的完整檢定列於附錄~\ref{sec:stats-full}。AUC 為主要報告指標，但其改善不自動代表指定門檻下的 Recall 也同時改善；三項指標須分別解讀。

\subsection{不同測試年度的表現}

\input{generated/chapter-results-20261027/ensemble-yearly.tex}

表~\ref{tab:results-ensemble-yearly} 分開呈現四個測試年度。以 V-static 為例，2015 年與 2018 年 AUC 為 0.8022 與 0.8556，F1-score 卻由 0.2949 降至 0.1304。這說明年度排序區辨能力與固定門檻下的分類表現可能呈現不同變化，不能只由 pooled AUC 判斷各年度的辨識情形。

測試正類紀錄數由 2015 年的 111 筆降至 2018 年的 36 筆，類別比例與樣本組成也隨年度不同。年度結果維持原驗證門檻，沒有依測試年度重新校正；本表不將四個年度當成四次獨立實驗，也未據此宣稱變化由特定景氣或漂移因素造成。

\section{實驗二：FS 集成比較}\label{sec:results-experiment2}

接著將八種 FS 條件，與相同種子、測試公司及集成方法的 no-FS 結果配對。兩種保留設定分別使用 9 與 14 個特徵，精簡率為 50.00\% 與 22.22\%。各條件都重新訓練模型，並依相同的 Calibration 規則選擇靜態成員與門檻。

\subsection{特徵條件的完整測試結果}

表~\ref{tab:results-fs-auc}、表~\ref{tab:results-fs-f1} 與表~\ref{tab:results-fs-recall} 分別列出八個 FS 條件及 no-FS 對照。每一欄使用同一種集成規則，每一條件以十個種子彙整，不把不同方法的分數混成單一平均。

\input{generated/chapter-results-20261027/fs-auc-absolute.tex}
\input{generated/chapter-results-20261027/fs-f1-absolute.tex}
\input{generated/chapter-results-20261027/fs-recall-absolute.tex}

以 Equal6 為例，no-FS、MI 50\% 與 MI 80\% 的 AUC 為 0.7815、0.7644 與 0.7822，顯示大幅刪減與保留較多特徵的結果不同。V-static 在 SHAP 50\% 下的 Recall 為 0.3502，高於 no-FS 的 0.2481，但 F1-score 為 0.1867，低於 no-FS 的 0.2355。因而不能把少數類涵蓋率的增加直接視為所有指標均改善。

\subsection{與同方法 no-FS 的配對比較}

\input{generated/statistics-20261027T051346477136Z/experiment2-auc-differences.tex}

表~\ref{tab:stats-exp2-auc} 呈現各 FS 條件相對同方法 no-FS 的 AUC 差值。完整的 48 組 AUC 比較，以及對應的 F1-score、Recall 檢定與信賴區間，均列於附錄~\ref{sec:stats-full}。差值的方向與顯著性分開判讀，未達顯著不代表已證明方法等效。

在本次整體校正下，沒有 FS 條件呈現顯著的 AUC 增加，F1-score 的全部 48 項比較亦未達顯著。CART 50\% 搭配 KNORA-E 的 AUC 差值為 \(-0.0435\)，其 \(p_H=0.03900\)，方向為下降；該值接近顯著門檻且受有限次 Monte Carlo 誤差影響，因此不將此單一門檻判斷視為穩健結論。

Recall 共有九項比較通過 Holm 校正，其中 V-static 在 MI 80\%、SHAP 50\%、SHAP 80\% 及 RFE 80\% 下增加，其餘五項為 DCS 方法下的下降。例如，SHAP 50\% 的 V-static Recall 增加 0.1021，\(p_H=0.00632\)；CART 50\% 的 DCS-LCA Recall 降低 0.1847，\(p_H=0.00330\)。此差異顯示 FS 的影響與集成方法有關，不能僅因部分 Recall 改善，就推論 FS 全面提升區辨能力或精確率與召回率的平衡。

對 V-static 而言，FS 後的靜態成員與門檻可依 Calibration 重新決定，其差異反映整體程序的影響。對固定方法而言，F1-score 與 Recall 仍可能受到 FS 後的模型機率與驗證門檻變化共同影響，因此不將 Recall 的變化單獨解釋為特徵數減少的直接效果。

\section{統計推論範圍與限制}

本章檢定屬於既有結果的事後探索分析，完整保留 165 項比較，避免只報告增加或達顯著者。公司群聚重抽考量同一公司的重複測試紀錄，但使用已固定的模型與驗證決策，不涵蓋重新取得訓練資料、重新選模與其他年度環境的全部不確定性。不同公司也可能受到共同產業或景氣衝擊，這些相依性並未另行建模。

本章使用公開資料中的公司層級狀態，該標籤在同一公司所有年度保持相同，且公司可能跨擬合與測試期間出現。因此，結論限於本資料與既定時間分區下的回溯性公司狀態分類。統計差異不等同模型已具備提前一年預測能力，也不能據此宣稱對全新公司或不同市場具有普遍優越性。第三章的欄位表已依實際發布版修正，但模型使用的財務欄位、數值與預測結果保持原狀。
\endgroup
